% ============================================================================
% 2.4 平行数组：字符串的索引存储
% ============================================================================
\subsection{为什么需要索引存储}

每一次循环，用户都会输入一行内容。如果只是把它们零散地放到堆里，
时间久了就找不到了。我们需要一种\textbf{按序号}存取的结构：
存第 0 个、第 1 个、第 2 个……，然后按同样的序号取出来打印。

为此我们用两个平行数组：一个存每个字符串的\textbf{指针}（64 位 / 8 字节），
一个存每个字符串的\textbf{长度}（32 位 / 4 字节）。配合一个计数器记录
存了多少条。

\subsection{string\_array.s：数据区定义}

\begin{lstlisting}[style=asm,caption=string\_array.s — 数据区定义]
.data
.align 4

.equ MAX_STRINGS, 32             ; 最多存 32 条字符串

string_count: .word 0            ; 计数器，已存入多少条
string_ptrs: .space MAX_STRINGS * 8  ; 指针数组，每项 8 字节
string_lens: .space MAX_STRINGS * 4  ; 长度数组，每项 4 字节
\end{lstlisting}

从内存地址看，它的布局是：

\begin{center}
\begin{tabular}{|l|l|l|l|}
\hline
string\_count & 4 字节 & 存放当前已存入条数 & \\
\hline
string\_ptrs[0] & 8 字节 & 第 0 条的地址 & \\
\hline
string\_ptrs[1] & 8 字节 & 第 1 条的地址 & \\
\hline
\(\vdots\) & & & \\
\hline
string\_ptrs[31] & 8 字节 & 第 31 条的地址 & \\
\hline
string\_lens[0] & 4 字节 & 第 0 条的长度 & \\
\hline
\(\vdots\) & & & \\
\hline
\end{tabular}
\end{center}

\subsection{store\_string\_ptr\_len：写入一对指针和长度}

\begin{lstlisting}[style=asm,caption=string\_array.s — store\_string\_ptr\_len]
.text
.global store_string_ptr_len
.global store_ptr_done
.global get_string_ptr
.global get_string_len
.global print_index_strings
.extern sys_write

store_string_ptr_len:
    stp x29, x30, [sp, #-16]!
    mov x29, sp

    ; 检查是否超过上限
    cmp x1, #MAX_STRINGS
    b.hs store_ptr_done

    ; 存指针：string_ptrs[x1 * 8] = x0
    adrp x4, string_ptrs@PAGE
    add  x4, x4, string_ptrs@PAGEOFF
    str  x0, [x4, x1, lsl #3]   ; x1 * 8 作为偏移

    ; 存长度：string_lens[x1 * 4] = x2
    adrp x3, string_lens@PAGE
    add  x3, x3, string_lens@PAGEOFF
    lsl  x5, x1, #2
    str  w2, [x3, x5]

store_ptr_done:
    ldp x29, x30, [sp], #16
    ret
\end{lstlisting}

\paragraph{调用约定}
调用者需要把三个参数放进来：
\begin{itemize}
    \item \texttt{x0} — 字符串在内存中的起始地址（就是 \texttt{copy\_string} 的返回值）
    \item \texttt{x1} — 序号（第几条）
    \item \texttt{x2} — 字符串长度（字节数）
\end{itemize}

注意 ARM 的指令 \texttt{str x0, [x4, x1, lsl \#3]} 可以直接
在一个内存操作里完成"基址 + 变址×8"的计算——因为指针是 8 字节，
正好用 \texttt{lsl \#3}（左移 3 位 = ×8）。长度是 4 字节，
所以用 \texttt{lsl \#2}（左移 2 位 = ×4），但 AArch64 不允许
\texttt{str} 的寻址直接放 \texttt{lsl \#2}，所以先算到 \texttt{x5} 再用。

\subsection{print\_index\_strings：按序号把字符串打印出来}

写入之后我们要能读回来打印。读的逻辑是写入的反向：

\begin{lstlisting}[style=asm,caption=string\_array.s — print\_index\_strings]
print_index_strings:
    stp x29, x30, [sp, #-16]!
    mov x29, sp

    cmp x0, #MAX_STRINGS
    b.hs print_index_strings_done

    ; 取指针：x1 = string_ptrs[x0 * 8]
    adrp x3, string_ptrs@PAGE
    add  x3, x3, string_ptrs@PAGEOFF
    ldr  x1, [x3, x0, lsl #3]

    ; 取长度：x2 = string_lens[x0 * 4]
    adrp x3, string_lens@PAGE
    add  x3, x3, string_lens@PAGEOFF
    lsl  x4, x0, #2
    ldr  w2, [x3, x4]

    ; 如果指针或长度为 0，说明这个位置从未被写入过
    cbz x1, print_index_strings_done
    cbz x2, print_index_strings_done

    ; 调用 sys_write(1, ptr, len) 打印
    mov x0, #1
    stp x1, x2, [sp, #-16]!    ; 压栈保护 x1,x2
    bl  sys_write
    ldp x1, x2, [sp], #16

print_index_strings_done:
    ldp x29, x30, [sp], #16
    ret
\end{lstlisting}

调用者只需要把序号（0、1、2…）放 \texttt{x0}，\texttt{bl print\_index\_strings}
就能把那一行打印到标准输出。如果那个序号的位置从未写入过，
\texttt{string\_ptrs[i]} 里是 0（\texttt{.space} 初始化为 0），
就会被 \texttt{cbz} 抓住，直接退出——不会因为野指针而崩溃。

\paragraph{为什么用两个数组而不是一个结构体数组？}
如果用结构体（每条存一个指针 + 一个长度，12 字节对齐到 16 字节），
每次写入都要写一条连续的 16 字节，而读的时候不需要长度的话也要
把整条读出来。分开两个数组时：
\begin{itemize}
    \item 写入：两条 store 指令
    \item 只需要指针时：只读 8 字节
    \item 只需要长度时：只读 4 字节
    \item 按序号索引时，不需要算结构体对齐的 padding
\end{itemize}
对于这个小程序来说两种方式都可以，但平行数组的写法更直观，
循环内每次只做一件小事。
